\documentclass[../thesis.tex]{subfiles}

\begin{document}
	Тут будет проверка решения на соответствие требованиям.

	Для практическое верификации полученного решения были разработаны два сервиса, использующие все
	описанные ранее концепции. Данные сервисы могут быть использованы в качестве шаблона для инициализации
	последующих проектов. Подробнее с техническим обоснованием и результатами можно ознакомиться в
	Приложениях \hyperref[apx:2]{Б} и \hyperref[apx:3]{В}.

	\section{Интеграционный план}

	Для непрерывной интеграции в рабочий процесс, предлагается следующий план:

	\begin{enumerate}
		\item Внедрение брокера;

		\item Внедрение платформы сообщений.
	\end{enumerate}

	\section{Слабые места}

	\begin{itemize}
		\item Производительность графического интерфейса.
	\end{itemize}

	\section{Дальнейшее развитие}

	\begin{itemize}
		\item Разработка генератора графического интерфейса по схемам конфигурационных файлов;

		\item Создание изолированных управляемых пространств разработки.
	\end{itemize}
\end{document}